\chapter{Implementation Manual: Every Function and Its Execution Order}
\label{chap:function-manual}
This manual explains every named function, class method and nested callback in the six executable source modules. It complements the mathematical chapters with the actual implementation order, arguments, state changes, return values and interactions. Constants, module initialization, lambdas and main guards are explained separately. The appendix contains the complete source snapshot, including imports, tables, comments and entry points, so this document can be read without opening an external Python file.

\section{Reading Python data and control flow}
\begin{itemize}
\item A byte is an integer from zero to 255; immutable \texttt{bytes} are used for external blocks and stored snapshots. \texttt{list} and nested row lists hold mutable working states. Reassigning an element of an AES matrix changes that matrix; replacing a local integer requires returning the replacement.
\item \texttt{\^{}} is XOR, \texttt{\&} is bitwise AND, and shifts move bits. \texttt{0xFF} keeps eight low bits and \texttt{0xFFFFFFFF} keeps thirty-two. These operators are not ordinary addition and multiplication.
\item \texttt{[a:b]} copies a slice; \texttt{[::-1]} reverses order. AES row rotation moves bytes inside one row; SM4's final reversal moves four entire words. These operations must not be confused.
\item \texttt{None} means no value/observer or no observation yet. A nonempty candidate set with multiple elements means ambiguity. Empty intersections are model errors, not successful recovery of zero-valued key bytes.
\item An observer is a function passed to another function. A closure captures values from its containing execution. A generator yields observations only when the solver asks for them. These constructs connect the existing engines to the experiments without copying the encryption algorithms.
\item Each \texttt{if \_\allowbreak{}\_\allowbreak{}name\_\allowbreak{}\_\allowbreak{} == "\_\allowbreak{}\_\allowbreak{}main\_\allowbreak{}\_\allowbreak{}"} block runs only on direct execution. Importing a module defines its helpers and constants without automatically starting its terminal demonstration or opening a Tk window.
\end{itemize}

\section{Module-level configuration and initialization}
\begin{description}
\item[\texttt{aes.py}] \texttt{KEY} and \texttt{PLAINTEXT} are editable demonstration strings, not implicit inputs to the block functions. \texttt{BLOCK\_\allowbreak{}SIZE=16} and \texttt{N\_\allowbreak{}ROUNDS=10} select AES-128. \texttt{SBOX} contains the standard substitution table; \texttt{INV\_\allowbreak{}SBOX} is constructed by reversing its mapping. \texttt{RCON} uses an unused index zero so schedule rounds index directly. \texttt{MIX\_\allowbreak{}MATRIX} and \texttt{INV\_\allowbreak{}MIX\_\allowbreak{}MATRIX} contain the field coefficients for forward/inverse column mixing.
\item[\texttt{sm4.py}] The example strings use the official vector. \texttt{WORD\_\allowbreak{}MASK} bounds words to thirty-two bits; \texttt{BLOCK\_\allowbreak{}SIZE} and \texttt{N\_\allowbreak{}ROUNDS} select sixteen bytes and thirty-two rounds. The SM4 \texttt{SBOX} is different from AES. Four \texttt{FK} words initialize the key schedule; thirty-two \texttt{CK} words enter the schedule recurrence. Their byte construction is documented next to the constants in the source.
\item[\texttt{piret.py}] Imports expose the real cipher engines, Cartesian products for correlated quartet candidates, a cached inverse binary matrix, random fault sequences and terminal formatting. The legacy key/plaintext/fault globals are for \texttt{relation\_\allowbreak{}demo}; the main recovery path reads the cipher configuration or explicit parameters. The inverse AES S-box and mixing matrix are shared mathematical data, not a secret key.
\item[\texttt{giraud.py}] \texttt{SINGLE\_\allowbreak{}BITS} enumerates the eight possible one-bit byte masks. \texttt{AES\_\allowbreak{}INVERSE\_\allowbreak{}SBOX} is public algorithm data. The inverse SM4 data diffusion function is imported from the common attack module. Mask knowledge is an explicit argument, not inferred from the simulator's hidden state.
\item[\texttt{run.py}] Imports the target engines and common attack mathematics. \texttt{LessonTerminal} extends \texttt{piret.Terminal} with flat byte comparisons; its inherited color and pause behavior remains presentation only.
\item[\texttt{gui.py}] Imports JSON export, queue/thread support and Tk widgets, then the cipher/attack modules. \texttt{AES\_\allowbreak{}KEY}, \texttt{AES\_\allowbreak{}INPUT} and \texttt{SM4\_\allowbreak{}VECTOR} define examples. Color constants style the window. Its extra inverse-S-box table is module-level teaching data; candidate solving still delegates to the attack modules. The main guard creates the Tk root, constructs \texttt{CryptoLab} and enters \texttt{mainloop}.
\end{description}

\section{End-to-end trace: ordinary AES and SM4}
\begin{enumerate}
\item In AES, \texttt{aes\_\allowbreak{}encrypt(P,K)} first calls \texttt{bytes\_\allowbreak{}to\_\allowbreak{}state(P)} and \texttt{key\_\allowbreak{}expansion(K)}. The initial state byte at row one, column zero is \hex{11}; its matching first-key byte is \hex{01}, so the initial XOR produces \hex{10}.
\item \texttt{\_\allowbreak{}operation} applies the initial key and each subsequent transformation. At round one, \texttt{sub\_\allowbreak{}bytes} replaces \hex{10} by \hex{CA}; \texttt{shift\_\allowbreak{}rows} moves it to another column; \texttt{mix\_\allowbreak{}columns} combines column bytes using \texttt{gf\_\allowbreak{}mul}. \texttt{add\_\allowbreak{}round\_\allowbreak{}key} then applies K1.
\item The AES loop repeats until round ten, whose conditional skips MixColumns. \texttt{state\_\allowbreak{}to\_\allowbreak{}bytes} returns \hex{69C4E0D8\allowbreak{}6A7B0430\allowbreak{}D8CDB780\allowbreak{}70B4C55A} for the official vector. Decryption starts with K10 and reverses the transformations to recover P.
\item In SM4, \texttt{sm4\_\allowbreak{}encrypt\_\allowbreak{}block(P,K)} delegates to \texttt{sm4\_\allowbreak{}process\_\allowbreak{}block}. It unpacks P into X0 through X3 and calls \texttt{key\_\allowbreak{}expansion}. The first encryption subkey is \hex{F12186F9}.
\item \texttt{round} forms A0=\hex{F002C39E}, applies \texttt{sbox} to obtain \hex{18E992B1}, and \texttt{linear} to obtain \hex{26D99622}. XOR with X0=\hex{01234567} gives X4=\hex{27FAD345}. The returned window is X1,X2,X3,X4.
\item After thirty-two updates, SM4 serializes X35,X34,X33,X32, producing \hex{681EDF34\allowbreak{}D206965E\allowbreak{}86B3E94F\allowbreak{}536E4246}. The decryption wrapper uses the same recurrence with the subkeys in reverse order and checks the original block is recovered.
\end{enumerate}

\section{End-to-end trace: the four fault-recovery paths}
\begin{enumerate}
\item \textbf{AES round-nine Piret:} \texttt{pdf\_\allowbreak{}aes\_\allowbreak{}lesson} or \texttt{run.aes\_\allowbreak{}demo} asks \texttt{encrypt\_\allowbreak{}with\_\allowbreak{}trace} for a byte fault. Its observer mutates the real AES state before MixColumns. \texttt{key\_\allowbreak{}candidates} receives only the output pair and column; intersections recover a quartet. After four unique groups, \texttt{reverse\_\allowbreak{}key\_\allowbreak{}schedule} reconstructs K0. The direct \texttt{piret.py} command instead uses the documented round-eight generator and \texttt{recover\_\allowbreak{}last\_\allowbreak{}key} to process all groups together.
\item \textbf{AES Giraud:} \texttt{aes\_\allowbreak{}fault\_\allowbreak{}encrypt} changes one bit before final SubBytes. \texttt{aes\_\allowbreak{}candidates} identifies the single changed output byte and tests the inverse-S-box equation for every byte guess. \texttt{intersect\_\allowbreak{}candidates} retains shared survivors. Each unique output byte fills its K10 position; reversal and re-encryption follow only after all sixteen positions are unique.
\item \textbf{SM4 byte DFA:} \texttt{recover\_\allowbreak{}sm4\_\allowbreak{}master\_\allowbreak{}key} asks an output-only oracle for a selected round/byte/mask. \texttt{peel\_\allowbreak{}sm4} removes already recovered later rounds, and \texttt{sm4\_\allowbreak{}byte\_\allowbreak{}candidates} removes L from the difference. Intersections identify a subkey byte. All sixteen subkey bytes across rounds 31--28 must be unique before \texttt{reverse\_\allowbreak{}sm4} returns the master key.
\item \textbf{SM4 Giraud-inspired:} the same outer recovery loop runs, but the wrapper supplies \texttt{sm4\_\allowbreak{}candidates} through \texttt{candidate\_\allowbreak{}filter}. Unknown-bit mode tests all eight bit masks and never forwards the actual simulator mask. Controlled target access at all four rounds is still required; changing the filter does not remove that assumption.
\end{enumerate}
These paths share data structures and target engines, but their observation requirements and allowed fault sets differ. A teaching observer is allowed to show an intermediate value; the recovery API is not automatically entitled to receive it.

\section{How to use the function reference}
Every entry below gives the actual callable signature, its source line range in the appendix, and ordered implementation steps. Methods include \texttt{self}, the application or terminal instance on which they operate. Nested callbacks are described after their containing module. The appendix retains decorators, including the cached inverse-matrix decorator, as well as complete parent functions containing the callbacks. Signature wrapping in the typeset listing does not change the one-line argument convention in the source files. Code uses plain monospaced rendering to reduce compilation work; all source text is retained. The contents lists chapters and sections, while individual function headings remain in the body.
